\let\stop\empty
\documentclass{exam-zh}
\usepackage{siunitx}
\usepackage{tkz-euclide}
\usepackage{diagbox}  
\newenvironment{answer}{\par\vspace{0.5em}\noindent\textbf{【答案】}\par}{\par\vspace{1em}}

\newenvironment{solutions}{\par\vspace{0.5em}\noindent\textbf{【解析】}\par}{\par\vspace{1em}}\usepackage{lastpage}
\usepackage{caption}
\examsetup{
  page/size=a4paper,
  paren/show-paren=true,
  paren/show-answer=true,
  fillin/show-answer=false,
  font      = times,
  math-font = xits,
  fillin/no-answer-type=none,
  %solution/show-solution=show-stay,
  %solution/label-indentation=false
  }
%\newcommand{\customcurve}[1][1]{%
  %\tikz[baseline,scale=#1,line width=0.8pt,line cap=round]{%
    %\draw (0.15,0.1) ellipse (0.15 and 0.1);
    %\draw (0,0.1) -- (-0.12 , -0.1);
    %\draw (0,0.1) -- (-0.12 , 0.3);
  %}%
%}
\newcommand{\customcurve}[1][1]{%
  \tikz[baseline,yshift=3,scale=0.05,thick,line width=0.8pt,
        line cap=round, domain=-1:2.832, samples=300]{
    \draw plot (\x,{sqrt(16/(\x+2)^2 - (\x-2)^2)});
    \draw plot (\x,{-sqrt(16/(\x+2)^2 - (\x-2)^2)});
    }
}

\everymath{\displaystyle}

\title{2026年06月27日 试卷类模板组卷1}

\subject{数学}



\begin{document}

% \information{
%   姓名\underline{\hspace{6em}},
%   座位号\underline{\hspace{15em}}
% }
% \warning{（在此卷上答题无效）}

\secret

\maketitle
\begin{center}
    本试卷共 \pageref{LastPage} 页，19 题。全卷满分 150 分。考试用时 120 分钟。
\end{center}


\begin{notice}
  \item 答题前，先将自己的姓名、准考证号、考场号、座位号填写在试卷和答题卡上，
    并将准考证号条形码粘贴在答题卡上的指定位置。
  \item 选择题的作答：每小题选出答案后，用 2B 铅笔把答题卡上对应题目的答案标号涂黑。
    写在试卷、草稿纸和答题卡上的非答题区域均无效。
  \item 填空题和解答题的作答：用黑色签字笔直接答在答题卡上对应的答题区域内。
    写在试卷、草稿纸和答题卡上的非答题区域均无效。
  \item 考试结束后，请将本试卷和答题卡一并上交。
\end{notice}



% 1.
% === Begin Label Data ===
% ID: 837
% 难度星级: 1
% 标签: 空间图形，立体几何，命题与逻辑，概念辨析，空间几何体
% 备注: 
% 组卷引用次数: 3
% === End  Label Data ===

\begin{question}
```
下列说法中正确的是 (\hspace{1cm})
\begin{choices}
\item 长方体是四棱柱, 直四棱柱是长方体
\item 有 2 个面平行, 其余各面都是梯形的几何体是棱台
\item 各侧面都是正方形的四棱柱一定是正方体
\item 棱柱的侧棱相等, 侧面都是平行四边形
```
\end{choices}
\end{question}

\begin{answer}
$D$
\end{answer}

\begin{solutions}
选项 $A$ ：因为直四棱柱的底面可以是任意四边形，所以不一定是长方体。选项 $B$ ：因为棱台的侧棱延长线必须交于同一点，所以仅满足面平行且侧面为梯形不能判定为棱台。选项 $C$ ：因为底面可以为菱形，所以各侧面为正方形的四棱柱不一定是正方体。选项 $D$ ：因为棱柱的侧棱互相平行且长度相等，所以侧面必然均为平行四边形。综上所述，正确选项为 \boxed{D}。
\end{solutions}

% 2.
% === Begin Label Data ===
% ID: 840
% 难度星级: 2.5
% 标签: 向量运算，三角形的四心，坐标系法
% 备注: 
% 组卷引用次数: 3
% === End  Label Data ===

\begin{question}
```
(多选) $\triangle ABC$ 中, $AB=AC=5, BC=6$. 点 $P$ 满足 $\overrightarrow{AP}=x \overrightarrow{AB}+y \overrightarrow{AC}$, 设 $\lambda=x+y$, 则 (\hspace{1cm})
\begin{choices}
\item 若 $P$ 为 $\triangle ABC$ 的重心, 则 $\lambda=\frac{2}{3}$
\item 若 $P$ 为 $\triangle ABC$ 的内心, 则 $\lambda=\frac{5}{8}$
\item 若 $P$ 为 $\triangle ABC$ 的垂心, 则 $\lambda=\frac{7}{16}$
\item 若 $P$ 为 $\triangle ABC$ 的外心, 则 $\lambda=\frac{5}{6}$
```
\end{choices}
\end{question}

\begin{answer}
$ABC$
\end{answer}

\begin{solutions}
建立平面直角坐标系，以 $BC$ 中点为原点，$BC$ 所在直线为 $x$ 轴，$BC$ 边上的高为 $y$ 轴。
则 $A(0,4), B(-3,0), C(3,0)$.
$$
\overrightarrow{AB}=(-3,-4), \quad \overrightarrow{AC}=(3,-4)
$$
设 $P(x_P,y_P)$,由 $\overrightarrow{AP}=x \overrightarrow{AB}+y \overrightarrow{AC}$ 比较纵坐标得：
$$
y_P-4=-4(x+y) \Rightarrow \lambda=x+y=\frac{4-y_P}{4}
$$
该式表明 $\lambda$ 仅与点 $P$ 的纵坐标有关。

（$1$）若 $P$ 为重心，则 $y_P=\displaystyle\frac{4+0+0}{3}=\displaystyle\frac{4}{3}.$
代入得 $\lambda=\displaystyle\frac{4-\frac{4}{3}}{4}=\displaystyle\frac{2}{3}.$ 所以选项 $A$ 正确。

（$2$）若 $P$ 为内心，则 $y_P=\displaystyle\frac{6\times 4+5\times 0+5\times 0}{6+5+5}=\displaystyle\frac{3}{2}.$
代入得 $\lambda=\displaystyle\frac{4-\frac{3}{2}}{4}=\displaystyle\frac{5}{8}.$ 所以选项 $B$ 正确。

（$3$）若 $P$ 为垂心，由对称性知垂心在 $y$ 轴上。
$AC$ 边斜率为 $-\displaystyle\frac{4}{3}$,故 $AC$ 边上高线方程为 $y=\displaystyle\frac{3}{4}(x+3).$
令 $x=0$ 得 $y_P=\displaystyle\frac{9}{4},$ 代入得 $\lambda=\displaystyle\frac{4-\frac{9}{4}}{4}=\displaystyle\frac{7}{16}.$ 所以选项 $C$ 正确。

（$4$）若 $P$ 为外心，设 $P(0,y_P),$ 由 $PA=PB$ 得 $(4-y_P)^2=3^2+y_P^2.$
解得 $y_P=\displaystyle\frac{7}{8},$ 代入得 $\lambda=\displaystyle\frac{4-\frac{7}{8}}{4}=\displaystyle\frac{25}{32}\neq\displaystyle\frac{5}{6}.$ 所以选项 $D$ 错误。

综上所述，本题正确选项为 $ABC.$ \boxed{ABC}
\end{solutions}

% 3.
% === Begin Label Data ===
% ID: 846
% 难度星级: 2.0
% 标签: 平面向量，向量数量积，向量的模
% 备注: 
% 组卷引用次数: 3
% === End  Label Data ===

\begin{question}
已知向量 $\boldsymbol{a}$, $\boldsymbol{b}$ 满足 $|\boldsymbol{a}|=1$, $|\boldsymbol{a}+2\boldsymbol{b}|=2$, 且 $(\boldsymbol{b}-2\boldsymbol{a}) \perp \boldsymbol{b}$, 则 $|\boldsymbol{b}|=$ (\hspace{1cm})
\begin{choices}
\item $\frac{1}{2}$
\item $\frac{\sqrt{2}}{2}$
\item $\frac{\sqrt{3}}{2}$
\item $1$
\end{choices}
\end{question}

\begin{answer}
$B$
\end{answer}

\begin{solutions}
思路 1 代数法.

由$(b-2a) \perp b$可知$(b-2a) \cdot b=0$,故$a \cdot b=\frac{1}{2}b^2$.

再由$|a+2b|=2$可得$a^2+4a \cdot b+4b^2=4$,故$b^2=\frac{1}{6}(4-a^2)$.

代入$|a|=1$可得$b^2=\frac{1}{2}$,故$|b|=\frac{\sqrt{2}}{2}$,正确选项是 B.

思路 2 坐标法.

由已知可设$a=(1,0)$,$b=(x,y)$,则$a+2b=(2x+1,2y)$,$b-2a=(x-2,y)$. 由已知得

$$
\begin{cases} (2x+1)^2+4y^2=4, \\ x(x-2)+y^2=0. \end{cases}
$$

解得$x=\frac{1}{4}$,$y=\pm\frac{\sqrt{7}}{4}$. 故$|b|=\sqrt{x^2+y^2}=\frac{\sqrt{2}}{2}$,应该选 B.

思路 3 几何法.

由题意构建如下平面图形，其中
\[
\overrightarrow{AP}=\overrightarrow{CD}=\frac{1}{2}\overrightarrow{AE}=\boldsymbol{a},\quad
\overrightarrow{AB}=\overrightarrow{GA}=\frac{1}{2}\overrightarrow{AC}=\boldsymbol{b},
\]
则
\[
\overrightarrow{AF}=\overrightarrow{BE}=2\boldsymbol{a}-\boldsymbol{b},\quad
\overrightarrow{AD}=\boldsymbol{a}+2\boldsymbol{b}.
\]

设$x=|b|=AB$,则$\cos \angle PAQ=\frac{AB}{AE}=\frac{x}{2}$,从而$\cos \angle ACD=-\frac{x}{2}$. 在$\triangle ACD$中，由余弦定理可得$AD^2=AC^2+CD^2-2AC \cdot CD \cdot \cos \angle ACD$,故

$$
4=4x^2+1^2+2 \cdot 2x \cdot 1 \cdot \frac{x}{2},
$$

解得$|b|=x=\frac{\sqrt{2}}{2}$. 所以，正确选项是 B.
\end{solutions}

% 4.
% === Begin Label Data ===
% ID: 843
% 难度星级: 3.5
% 标签: 球的内接几何体，球的截面性质，几何体体积
% 备注: 
% 组卷引用次数: 3
% === End  Label Data ===

\begin{question}
已知 $A, B, C$ 是半径为 1 的球 $O$ 的球面上的三个点，且 $AC \perp BC, AC=BC=1$，则三棱锥 $O-ABC$ 的体积为 (\hspace{1cm})
\begin{choices}
\item $\frac{\sqrt{2}}{12}$
\item $\frac{\sqrt{3}}{12}$
\item $\frac{\sqrt{2}}{4}$
\item $\frac{\sqrt{3}}{4}$
\end{choices}
\end{question}

\begin{answer}
$A$
\end{answer}

\begin{solutions}
$$
AB = \sqrt{AC^2 + BC^2} = \sqrt{2}
$$

因为 $\triangle ABC$ 为直角三角形，所以其外接圆半径 $r = \frac{AB}{2} = \frac{\sqrt{2}}{2}$ .

设球心 $O$ 到平面 $ABC$ 的距离为 $d$,则

$$
d = \sqrt{R^2 - r^2} = \sqrt{1 - \frac{1}{2}} = \frac{\sqrt{2}}{2}
$$

底面 $\triangle ABC$ 的面积

$$
S_{\triangle ABC} = \frac{1}{2} AC \cdot BC = \frac{1}{2}
$$

三棱锥 $O-ABC$ 的体积

$$
V = \frac{1}{3} S_{\triangle ABC} \cdot d = \frac{1}{3} \times \frac{1}{2} \times \frac{\sqrt{2}}{2} = \frac{\sqrt{2}}{12}
$$

故选 \boxed{A} .
\end{solutions}

% 5.
% === Begin Label Data ===
% ID: 842
% 难度星级: 3.5
% 标签: 圆锥侧面积，体积计算，比例关系
% 备注: 
% 组卷引用次数: 3
% === End  Label Data ===

\begin{question}
```
甲、乙两个圆锥的母线长相等，侧面展开图的圆心角之和为 $2 \pi$，侧面积分别为 $S_{\text{甲}}$ 和 $S_{\text{乙}}$，体积分别为 $V_{\text{甲}}$ 和 $V_{\text{乙}}$. 若 $\frac{S_{\text{甲}}}{S_{\text{乙}}}=2$，则 $\frac{V_{\text{甲}}}{V_{\text{乙}}}=$ (\hspace{1cm})
\begin{choices}
\item $\sqrt{5}$
\item $2 \sqrt{2}$
\item $\sqrt{10}$
\item $\frac{5 \sqrt{10}}{4}$
```
\end{choices}
\end{question}

\begin{answer}
$C$
\end{answer}

\begin{solutions}
设母线长为$l$,甲圆锥底面半径为$r_1$,乙圆锥底面圆半径为$r_2$,则
$$
\frac{S_{\text{甲}}}{S_{\text{乙}}} = \frac{\pi r_1 l}{\pi r_2 l} = \frac{r_1}{r_2} = 2,
$$
所以$r_1 = 2r_2$,又
$$
\frac{2\pi r_1}{l} + \frac{2\pi r_2}{l} = 2\pi,
$$
则$\frac{r_1 + r_2}{l} = 1$,所以$r_1 = \frac{2}{3}l, r_2 = \frac{1}{3}l$,所以甲圆锥的高$h_1 = \sqrt{l^2 - \frac{4}{9}l^2} = \frac{\sqrt{5}}{3}l$,乙圆锥的高$h_2 = \sqrt{l^2 - \frac{1}{9}l^2} = \frac{2\sqrt{2}}{3}l$,所以
$$
\frac{V_{\text{甲}}}{V_{\text{乙}}} = \frac{\frac{1}{3}\pi r_1^2 h_1}{\frac{1}{3}\pi r_2^2 h_2} = \frac{\frac{4}{9}l^2 \times \frac{\sqrt{5}}{3}l}{\frac{1}{9}l^2 \times \frac{2\sqrt{2}}{3}l} = \sqrt{10}.
$$
故选：C.
\end{solutions}

% 6.
% === Begin Label Data ===
% ID: 838
% 难度星级: 5
% 标签: 概率计算，条件概率，分类讨论
% 备注: 
% 组卷引用次数: 3
% === End  Label Data ===

\begin{question}{2023}{G}{新高考Ⅱ卷}{12}{概率}
```
（多选）在信道内传输 $0,1$ 信号，信号的传输相互独立，发送 0 时，收到 1 的概率为 $\alpha(0<\alpha<1)$；发送 1 时，收到 0 的概率为 $\beta(0<\beta<1)$ 收到 1 的概率为 $1-\beta$.考虑两种传输方案：单次传输和三次传输。单次传输是指每个信号只发送 1 次；三次传输是指每个信号重复发送 3 次。收到的信号需要译码，译码规则如下：单次传输时，收到的信号即为译码；三次传输时，收到的信号中出现次数多的即为译码（例如，若依次收到 $1,0,1$，则译码为 1）。 (\hspace{1cm})
\begin{choices}
\item 采用单次传输方案，若依次发送 $1,0,1$，则依次收到 $1,0,1$ 的概率为 $(1-\alpha)(1-\beta)^{2}$
\item 采用三次传输方案，若发送 1，则依次收到 $1,0,1$ 的概率为 $\beta(1-\beta)^{2}$
\item 采用三次传输方案，若发送 1，则译码为 1 的概率为 $\beta(1-\beta)^{2}+(1-\beta)^{3}$
\item 当 $0<\alpha<0.5$，若发送 0，则采用三次传输方案译码为 0 的概率大于采用单次传输方案译码为 0 的概率
\end{choices}
\end{question}

\begin{answer}
$ABD$
\end{answer}

\begin{solutions}
发送 $1$ 时收到 $1$ 的概率为 $1-\beta$,发送 $0$ 时收到 $0$ 的概率为 $1-\alpha$.
对于选项 $A$,依次发送 $1,0,1$ 收到 $1,0,1$ 的概率为 $(1-\beta)(1-\alpha)(1-\beta)=(1-\alpha)(1-\beta)^{2}$,故 $A$ 正确.
对于选项 $B$,发送 $1$ 三次依次收到 $1,0,1$ 的概率为 $(1-\beta)\beta(1-\beta)=\beta(1-\beta)^{2}$,故 $B$ 正确.
对于选项 $C$,发送 $1$ 译码为 $1$ 需收到两个或三个 $1$,概率为 $C_{3}^{2}(1-\beta)^{2}\beta+(1-\beta)^{3}=3\beta(1-\beta)^{2}+(1-\beta)^{3}$,故 $C$ 错误.
对于选项 $D$,单次传输译码为 $0$ 的概率为 $1-\alpha$. 三次传输译码为 $0$ 的概率为 $C_{3}^{2}(1-\alpha)^{2}\alpha+(1-\alpha)^{3}=3\alpha(1-\alpha)^{2}+(1-\alpha)^{3}$.
作差得 $[3\alpha(1-\alpha)^{2}+(1-\alpha)^{3}]-(1-\alpha)=\alpha(1-\alpha)(1-2\alpha)$.
当 $0<\alpha<0.5$ 时，$\alpha(1-\alpha)(1-2\alpha)>0$,故三次传输概率更大，$D$ 正确.
综上选 $ABD$.
\end{solutions}

% 7.
% === Begin Label Data ===
% ID: 844
% 难度星级: 3.0
% 标签: 正弦定理，余弦定理，三角恒等变换
% 备注: 
% 组卷引用次数: 3
% === End  Label Data ===

\begin{question}
记 $\triangle ABC$ 的内角 $A,B,C$ 的对边分别为 $a,b,c$, 已知 $\sin C\sin(A-B)=\sin B\sin(C-A)$.

（1） 证明: $2a^2=b^2+c^2$;

（2)）若 $a=5, \cos A=\frac{25}{31}$, 求 $\triangle ABC$ 的周长.
\end{question}

\begin{answer}
$（1）见解析
（2）14$
\end{answer}

\begin{solutions}
【分析】（1）利用两角差的正弦公式化简，再根据正弦定理和余弦定理化角为边，从而即可得证；

（2）根据（1）的结论结合余弦定理求出$bc$,从而可求得$b+c$,即可得解.

【小问1详解】

证明：因为$\sin C\sin(A-B)=\sin B\sin(C-A)$,

所以$\sin C\sin A\cos B-\sin C\sin B\cos A=\sin B\sin C\cos A-\sin B\sin A\cos C$,

所以$ac\cdot\frac{a^2+c^2-b^2}{2ac}-2bc\cdot\frac{b^2+c^2-a^2}{2bc}=-ab\cdot\frac{a^2+b^2-c^2}{2ab}$,

即$\frac{a^2+c^2-b^2}{2}-\left(b^2+c^2-a^2\right)=-\frac{a^2+b^2-c^2}{2}$,

所以$2a^2=b^2+c^2$;

【小问2详解】

解：因为$a=5,\cos A=\frac{25}{31}$,

由（1）得$b^2+c^2=50$,

由余弦定理可得$a^2=b^2+c^2-2bc\cos A$,

则$50-\frac{50}{31}bc=25$,

所以$bc=\frac{31}{2}$,

故$(b+c)^2=b^2+c^2+2bc=50+31=81$,

所以$b+c=9$,

所以$\triangle ABC$的周长为$a+b+c=14$.
\end{solutions}

% 8.
% === Begin Label Data ===
% ID: 835
% 难度星级: 5
% 标签: 函数性质，向量，三角函数
% 备注: 
% 组卷引用次数: 20
% === End  Label Data ===

\begin{question}{2026}{M}{5·3高中同步增分测评卷}{19}{向量，函数}
```
已知O为坐标原点, 对于函数 $f(x)=a \sin x+b \cos x$, 称向量 $\overrightarrow{O M}=(a, b)$ 为函数 $f(x)$ 的相伴特征向量, 同时称函数 $f(x)$ 为向量 $\overrightarrow{O M}$ 的相伴函数.
```
（1） 设函数 $g(x)=\sin \left(x+\frac{5 \pi}{6}\right)-\sin \left(\frac{3 \pi}{2}-x\right)$, 试求 $g(x)$ 的相伴特征向量 $\overrightarrow{O M}$;
```
（2） 记向量 $\overrightarrow{O N}=(1, \sqrt{3})$ 的相位函数为 $f(x)$, 或当 $f(x)=\frac{8}{5}$ 且 $x \in\left(-\frac{\pi}{2}, \frac{\pi}{6}\right)$ 时, $\sin x$ 的取值;```

（3） 已知 $A(-2,3), B(2, b), \overrightarrow{D T}=(-\sqrt{3}, 1)$ 为函数 $h(x)=\min \left(x-\frac{\pi}{4}\right)$ 的相位特征向量, $\varphi(x)=h\left(\frac{x}{2}-\frac{\pi}{3}\right)$, 在 $y=\varphi(x)$ 的图象上是否存在一点 $P$, 使得 $\overrightarrow{A P} \perp \overrightarrow{B P}$ ?若存在, 求出点 $P$ 的坐标; 若不存在, 请说明理由.
```
\end{question}

\begin{answer}
（1）$\boxed{\left(-\dfrac{\sqrt{3}}{2}, \dfrac{3}{2}\right)}$;（2）$\boxed{\dfrac{4-3\sqrt{3}}{10}}$;（3）存在，$\boxed{P(0,2)}$（当 $b=6$ 时）
\end{answer}

\begin{solutions}
（1）利用两角和与差的正弦公式展开化简。
$$
g(x)=\sin x\cos\frac{5\pi}{6}+\cos x\sin\frac{5\pi}{6}-(-\cos x)
$$
$$
=-\frac{\sqrt{3}}{2}\sin x+\frac{3}{2}\cos x
$$
因为相伴特征向量为 $(a,b)$,所以 $\overrightarrow{OM}=\left(-\frac{\sqrt{3}}{2}, \frac{3}{2}\right)$ .

（2）由定义得 $f(x)=\sin x+\sqrt{3}\cos x=2\sin\left(x+\frac{\pi}{3}\right)$ .
因为 $f(x)=\frac{8}{5}$,所以 $\sin\left(x+\frac{\pi}{3}\right)=\frac{4}{5}$ .
又 $x\in\left(-\frac{\pi}{2}, \frac{\pi}{6}\right)$,故 $x+\frac{\pi}{3}\in\left(-\frac{\pi}{6}, \frac{\pi}{2}\right)$ .
所以 $\cos\left(x+\frac{\pi}{3}\right)=\sqrt{1-\left(\frac{4}{5}\right)^2}=\frac{3}{5}$ .
因此 $\sin x=\sin\left[\left(x+\frac{\pi}{3}\right)-\frac{\pi}{3}\right]=\frac{4}{5}\times\frac{1}{2}-\frac{3}{5}\times\frac{\sqrt{3}}{2}=\frac{4-3\sqrt{3}}{10}$ .

（3）由 $\overrightarrow{DT}=(-\sqrt{3},1)$ 知 $h(x)=-\sqrt{3}\sin x+\cos x=2\cos\left(x+\frac{\pi}{3}\right)$ .
则 $\varphi(x)=h\left(\frac{x}{2}-\frac{\pi}{3}\right)=2\cos\left(\frac{x}{2}-\frac{\pi}{3}+\frac{\pi}{3}\right)=2\cos\frac{x}{2}$ .
设 $P(x_0,y_0)$ 在图象上，则 $y_0=2\cos\frac{x_0}{2}$ .
由 $\overrightarrow{AP}\perp\overrightarrow{BP}$ 得 $(x_0+2)(x_0-2)+(y_0-3)(y_0-b)=0$ .
取 $x_0=0$,则 $y_0=2$,代入得 $-4+(-1)(2-b)=0$,解得 $b=6$ .
当 $b=6$ 时，存在点 $P(0,2)$ 满足条件 .
若 $b\neq 6$,联立方程组分析可得无实数解 .
综上，当 $b=6$ 时存在点 $P(0,2)$;否则不存在 .
\end{solutions}

% 9.
% === Begin Label Data ===
% ID: 847
% 难度星级: 3
% 标签: 解三角形，三角恒等变换，余弦定理
% 备注: 
% 组卷引用次数: 3
% === End  Label Data ===

\begin{question}
已知 $\triangle ABC$ 的内角 $A, B, C$ 的对边分别为 $a, b, c$，且 $\sin(A+C)=8\sin^2\frac{B}{2}$.

（1）求 $\cos B$;

（2）若 $a+c=6$，$\triangle ABC$ 的面积为 2，求 $b$.
\end{question}

\begin{answer}
$\displaystyle\dfrac{15}{17}; 2$
\end{answer}

\begin{solutions}
（1）在 $\triangle ABC$ 中，因为 $A+B+C=\pi$,所以 $\sin(A+C)=\sin(\pi-B)=\sin B$.

原式化为 $\sin B=8\sin^2\frac{B}{2}$.

利用二倍角公式得 $2\sin\frac{B}{2}\cos\frac{B}{2}=8\sin^2\frac{B}{2}$.

因为 $B\in(0,\pi)$,所以 $\sin\frac{B}{2}\neq 0$,两边同除以 $2\sin\frac{B}{2}$ 得 $\cos\frac{B}{2}=4\sin\frac{B}{2}$.

平方得 $\cos^2\frac{B}{2}=16\sin^2\frac{B}{2}$.

结合 $\sin^2\frac{B}{2}+\cos^2\frac{B}{2}=1$,解得 $\sin^2\frac{B}{2}=\displaystyle\frac{1}{17}$.

由降幂公式 $\cos B=1-2\sin^2\frac{B}{2}=1-\displaystyle\frac{2}{17}=\displaystyle\frac{15}{17}$.

（2）因为 $\cos B=\displaystyle\frac{15}{17}$ 且 $B\in(0,\pi)$,所以 $\sin B=\sqrt{1-\cos^2 B}=\displaystyle\frac{8}{17}$.

由面积公式 $S_{\triangle ABC}=\displaystyle\frac{1}{2}ac\sin B=2$,得 $ac=\displaystyle\frac{4}{\sin B}=\displaystyle\frac{17}{2}$.

由余弦定理 $b^2=a^2+c^2-2ac\cos B=(a+c)^2-2ac(1+\cos B)$.

代入 $a+c=6$,$ac=\displaystyle\frac{17}{2}$,$\cos B=\displaystyle\frac{15}{17}$ 得：

$$
b^2=6^2-2\times\displaystyle\frac{17}{2}\times\left(1+\displaystyle\frac{15}{17}\right)=36-32=4.
$$

所以 $b=\boxed{2}$.
\end{solutions}

% 10.
% === Begin Label Data ===
% ID: 845
% 难度星级: 3.0
% 标签: 正弦定理，余弦定理，三角恒等变换，三角形的面积
% 备注: 
% 组卷引用次数: 3
% === End  Label Data ===

\begin{question}
记 $\triangle ABC$ 的内角 $A,B,C$ 的对边分别为 $a,b,c$，已知 $\sin C=\sqrt{2}\cos B, a^2+b^2-c^2=\sqrt{2}ab$.

（1）求角 B;

（2）若 $\triangle ABC$ 的面积为 $3+\sqrt{3}$，求 $c$.
\end{question}

\begin{answer}
（1）$\dfrac{\pi}{3}$;（2）$2\sqrt{2}$
\end{answer}

\begin{solutions}
(1)由已知及余弦定理可得$\cos C=\frac{a^{2}+b^{2}-c^{2}}{2ab}=\frac{\sqrt{2}}{2}$,

又$0<C<\pi$,故$C=\frac{\pi}{4}$.

由已知得$\cos B=\frac{\sin C}{\sqrt{2}}=\frac{1}{2}$,又$0<B<\pi$,故$B=\frac{\pi}{3}$.

(2)解法1

由(1)知$\sin A=\sin\left(\frac{\pi}{4}+\frac{\pi}{3}\right)=\frac{\sqrt{2}+\sqrt{6}}{4}$.

由(1)及正弦定理可得$b=\frac{\sqrt{6}}{2}c$,故$\triangle ABC$的面积

$S=\frac{1}{2}bc\sin A=\frac{3+\sqrt{3}}{8}c^{2}=3+\sqrt{3}$,

解得$c=2\sqrt{2}$.

解法2

由(1)可知$B=\frac{\pi}{3}$,$C=\frac{\pi}{4}$,故$\sin A=\sin(B+C)=\frac{\sqrt{2}+\sqrt{6}}{4}$.

由正弦定理$\frac{a}{\sin A}=\frac{c}{\sin C}$得$a=\frac{1+\sqrt{3}}{2}c$,故$\triangle ABC$的面积

$S=\frac{1}{2}ac\sin B=\frac{3+\sqrt{3}}{8}c^{2}=3+\sqrt{3}$,

解得$c=2\sqrt{2}$.

解法3

由(1)知$B=\frac{\pi}{3}$,$\sin C=\cos C=\frac{\sqrt{2}}{2}$,故$\sin A=\sin(B+C)=\frac{\sqrt{2}+\sqrt{6}}{4}$.

由正弦定理$\frac{a}{\sin A}=\frac{b}{\sin B}$得$a=\frac{3\sqrt{2}+\sqrt{6}}{6}b$,故$\triangle ABC$的面积

$S=\frac{1}{2}ab\sin C=\frac{3+\sqrt{3}}{12}b^{2}=3+\sqrt{3}$,

故$b=2\sqrt{3}$,$a=\sqrt{2}+\sqrt{6}$.

由已知可得$c^{2}=a^{2}+b^{2}-\sqrt{2}ab=8$,故$c=2\sqrt{2}$.
\end{solutions}

% 11.
% === Begin Label Data ===
% ID: 839
% 难度星级: 5.0
% 标签: 概率计算，分类讨论，独立事件
% 备注: 
% 组卷引用次数: 3
% === End  Label Data ===

\begin{question}
```
甲、乙、丙三位同学进行羽毛球比赛，约定赛制如下：  
累计负两场者被淘汰；比赛前抽签决定首先比赛的两人，另一人轮空；每场比赛的胜者与轮空者进行下一场比赛，负者下一场轮空，直至有一人被淘汰后，剩余的两人继续比赛，直到其中一人被淘汰，另一人最终获胜，比赛结束．  

经抽签，甲、乙首先比赛，丙轮空．设每场比赛双方获胜的概率都为$\frac{1}{2}$．  

（1）求甲连胜四场的概率；  
（2）求需要进行第五场比赛的概率；  
（3）求丙最终获胜的概率．
```
\end{question}

\begin{answer}
$\dfrac{1}{16}；\dfrac{3}{4}；\dfrac{37}{96}$
\end{answer}

\begin{solutions}
（1）甲连胜四场需连续赢得第 $1$, $2$, $3$, $4$ 场 .
因每场比赛双方获胜概率均为 $\displaystyle \frac{1}{2}$,且各场结果相互独立，
故所求概率为 $P_1 = \left(\displaystyle \frac{1}{2}\right)^4 = \displaystyle \frac{1}{16}$ .

（2）设 $A$ 为“需要进行第五场比赛” .
前四场比赛共有 $2^4 = 16$ 种等可能的胜负序列 .
比赛在第四场或之前结束，意味着有人在第四场前累计负两场 .
经状态转移分析，前四场未能决出最终胜者的序列恰有 $12$ 种 .
故 $P(A) = \displaystyle \frac{12}{16} = \displaystyle \frac{3}{4}$ .

（3）设 $B$ 为“丙最终获胜” .
丙首战轮空，从第二场起参赛 .
记 $x$ 表示“丙在场且己方负 $0$ 场时最终获胜的概率”，$y$ 表示“丙在场且己方负 $1$ 场时最终获胜的概率” .
根据赛制规则与全概率公式可列方程组：
$$
x = \displaystyle \frac{1}{2} + \displaystyle \frac{1}{2} y
$$
$$
y = \displaystyle \frac{1}{2} x
$$
解得 $x = \displaystyle \frac{2}{3}$,$y = \displaystyle \frac{1}{3}$ .
结合第一场甲乙赛果及后续对阵结构，
利用全概率公式综合计算得丙最终获胜的概率为 $P(B) = \displaystyle \frac{37}{96}$ .
\end{solutions}

\end{document}